\begin{abstract}
Optimizing eBPF programs is no longer a single compiler problem. A useful transformation may live in source code, in the LLVM BPF backend, in pre-load bytecode, in a live post-load rewrite, or in kernel/JIT extensions such as kinsn. Its value is decided only after the Linux verifier accepts the candidate, the application still behaves correctly, and the final JITed code improves the real workload. This makes eBPF optimization a natural target for agentic tuning, but current agent benchmarks do not exercise this verifier/JIT/workload loop.

We propose \sys, a benchmark for agents that tune existing eBPF programs under real feedback. A \sys task exposes an existing application, an action space over optimization layers, raw verifier/JIT/workload observations, and a composed oracle that rejects unsafe, broken, or merely noisy candidates. The first benchmark track focuses on bytecode, live ReJIT, and kinsn-backed actions over production-style eBPF applications, while keeping source-level and LLVM-backend transforms as explicit adapters. Historical results from our corpus show why this benchmark is needed: static pass policies are noisy, rewrite counts do not predict speedups, and tuned policies can move results in the right direction without solving the corpus. \sys turns these observations into a reproducible scoreboard for eBPF optimization agents.
\end{abstract}
